\section*{Capítulo \ref{ch:b2:conjugate-additions} --- Carbonilas conjugadas: Michael e Robinson}

\begin{solution}{exo:b2:conjugate-additions:1}
\begin{itemize}
  \item \ce{CH3MgBr}: sobretudo 1,2 (duro, irreversível), 1-metilciclo-hex-2-en-1-ol.
  \item \ce{(CH3)2CuLi}: 1,4, 3-metilciclo-hexanona.
  \item Tiofenol com base: 1,4 (o tiolato é mole), 3-(fenilsulfanil)ciclo-hexanona.
  \item \ce{LiAlH4}: 1,2, ciclo-hex-2-en-1-ol.
\end{itemize}
\end{solution}

\begin{solution}{exo:b2:conjugate-additions:2}
Doadores: pentano-2,4-diona, nitrometano, propanodioato de dietila (\ce{C-H} ácida). Aceptores:
propenonitrila, propenoato de metila, but-3-en-2-ona (\ce{C=C} conjugada com um grupo aceptor).
\end{solution}

\begin{solution}{exo:b2:conjugate-additions:3}
\ce{C4H9Br + 2Li -> C4H9Li + LiBr}, depois \ce{2C4H9Li + CuI -> (C4H9)2CuLi + LiI}, em éter a baixa
temperatura sob gás inerte.
\end{solution}

\begin{solution}{exo:b2:conjugate-additions:4}
(3-oxobutil)propanodioato de dietila, \ce{(EtOOC)2CHCH2CH2COCH3}. A hidrólise dá o ácido
propanodioico substituído, que perde \ce{CO2} sob aquecimento (\cref{prop:b2:enolates-aldol:decarboxylation}):
ácido 5-oxo-hexanoico, \ce{HOOCCH2CH2CH2COCH3}.
\end{solution}

\begin{solution}{exo:b2:conjugate-additions:5}
Dimetilcuprato de lítio em éter a baixa temperatura, depois tratamento aquoso: \omterm{def:b2:conjugate-additions:addition-modes}{adição conjugada}.
O brometo de metilmagnésio, mais duro, se adiciona sobretudo em 1,2 e dá o 1-metilciclo-hex-2-en-1-ol
como produto principal.
\end{solution}

\begin{solution}{exo:b2:conjugate-additions:6}
A \ce{Et3N} retira um pouco de \ce{EtSH}, dando \ce{EtS-}; o tiolato se adiciona ao carbono $\beta$, o
\omterm{def:b2:enolates-aldol:enolate}{enolato} formado toma o próton de outro \ce{EtSH}, o que regenera \ce{EtS-}. Produto:
4-(etilsulfanil)butan-2-ona. O enxofre é grande e polarizável, seu par não ligante alto em energia: um
\omterm{def:b2:frontier-orbitals:hard-soft}{nucleófilo mole}, sob \omterm{def:b2:frontier-orbitals:control}{controle orbital}, que portanto ataca onde a \omterm{def:b2:frontier-orbitals:frontier}{LUMO} é máxima; e a adição é
reversível, o que também favorece o aduto 1,4, mais estável.
\end{solution}

\begin{solution}{exo:b2:conjugate-additions:7}
O cianeto ataca mais depressa o carbono carbonílico, mais positivo: a cianoidrina é o produto cinético.
Sua formação é reversível; a temperatura mais alta, com tempo, o cianeto é liberado e se adiciona de
novo no carbono $\beta$, dando o aduto que conserva a ligação $\pi$ \ce{C=O}, mais forte: o produto
termodinâmico.
\end{solution}

\begin{solution}{exo:b2:conjugate-additions:8}
Com a 2-metilciclo-hexano-1,3-diona: a cetona de Wieland--Miescher (figura da \omterm{def:b2:conjugate-additions:robinson}{anelação de Robinson}).
Com a ciclo-hexanona: 4,4a,5,6,7,8-hexa-hidronaftalen-2(3\textit{H})-ona, uma enona bicíclica cuja
\ce{C=C} termina em um carbono de fusão, sem grupo metila angular.
\end{solution}

\begin{solution}{exo:b2:conjugate-additions:9}
\ce{CH3COCH2CH2CH2COCH3}: C2 e C6 são as carbonilas, cinco carbonos de C2 a C6 contados inclusive,
uma 1,5-dicetona. Corte C3--C4: doador, o \omterm{def:b2:enolates-aldol:enolate}{enolato} da propanona em C3 (melhor, o 3-oxobutanoato de etila,
cujo éster é removido depois por hidrólise e \omterm{def:b2:enolates-aldol:decarboxylation}{descarboxilação}); aceptor, a but-3-en-2-ona. Condições:
3-oxobutanoato de etila, but-3-en-2-ona, etóxido de sódio catalítico em etanol; depois ácido aquoso e
aquecimento.
\end{solution}

\begin{solution}{exo:b2:conjugate-additions:10}
A \omterm{def:b2:enolates-aldol:aldol}{aldólica} cria uma ligação \ce{C-C} mas transforma uma ligação $\pi$ \ce{C=O} em uma ligação $\sigma$
\ce{C-O}; o saldo é pequeno e duas moléculas tornam-se uma, de modo que $K$ é modesta e a \omterm{def:b2:enolates-aldol:aldol}{retroaldólica}
é fácil. A \omterm{def:b2:conjugate-additions:michael}{adição de Michael} cria uma ligação \ce{C-C} à custa de uma ligação $\pi$ \ce{C=C}, mais fraca,
e conserva as duas \ce{C=O}: ela é claramente favorável. A reação inversa exigiria que o \omterm{def:b2:enolates-aldol:enolate}{enolato} da cetona expulsasse o ânion propanodioato estabilizado; o
equilíbrio fica tão do lado do aduto que, com base catalítica, o aduto persiste.
\end{solution}

\begin{solution}{exo:b2:conjugate-additions:11}
\ce{C(COOEt)2(CH2CH2COOCH3)2}. Após a primeira \omterm{def:b2:conjugate-additions:michael}{adição de Michael}, o carbono central ainda tem um
hidrogênio entre dois ésteres, tão ácido quanto antes: a base o retira e segue-se uma segunda adição.
\end{solution}

\begin{solution}{exo:b2:conjugate-additions:12}
Pelo \omterm{def:b2:enolates-aldol:enolate}{enolato} mais substituído (do lado do carbono que leva a metila): a octalona com o grupo metila
no carbono de fusão, 4a-metil-4,4a,5,6,7,8-hexa-hidronaftalen-2(3\textit{H})-ona. Pelo \omterm{def:b2:enolates-aldol:enolate}{enolato} menos
substituído (do lado \ce{CH2}): o grupo metila termina em um carbono do anel vizinho da fusão. O
primeiro é favorecido em condições de equilíbrio (um alcóxido em seu álcool, ou a via da \omterm{def:b2:amines:imine}{enamina} com
uma \omterm{def:b2:amines:class}{amina secundária} e um ácido), que dão o \omterm{def:b2:enolates-aldol:kinetic-enolate}{enolato termodinâmico}, mais substituído
(\cref{prop:b2:enolates-aldol:enolate-control}).
\end{solution}

\begin{solution}{pb:b2:conjugate-additions:1}
\textbf{1.} Um anel de ciclo-hexano com \ce{C=O} em C1 e C3 e um grupo metila em C2; o hidrogênio em
C2 é o mais ácido.
\textbf{2.} Sua remoção dá um ânion cuja carga é deslocalizada sobre C2 e os dois oxigênios; na
ciclo-hexanona, um só oxigênio a compartilha, e a acidez é fraca demais para ser medida em água.
\textbf{3.} Carga em C2; \ce{C1=C2} com \ce{O^-} em C1; \ce{C2=C3} com \ce{O^-} em C3.
\textbf{4.} Sim: com o hidróxido, $K = 10^{14.0 - 5.3} = 10^{8.7}$ (a água como ácido conjugado, a partir
de $K_e$); com o etóxido, $10^{15.9-5.3} = 10^{10.6}$. Ambos são completos.
% ledger: ke:25C, pka:ethanol
\textbf{5.} \ce{C1(OH)=C2(CH3)}: C2 ainda leva um hidrogênio, que é tudo de que um \omterm{def:b2:enolates-aldol:tautomerism}{enol} precisa; o
grupo metila apenas substitui o segundo.
\textbf{6.} Mole: a carga está espalhada sobre três átomos e o ânion é estabilizado. Ele ataca o
carbono $\beta$ (\cref{prop:b2:conjugate-additions:regio}).
\textbf{7.} O C2 do \omterm{def:b2:enolates-aldol:enolate}{enolato} se liga à extremidade \ce{CH2} da but-3-en-2-ona; o \omterm{def:b2:enolates-aldol:enolate}{enolato} formado toma
um próton: 2-metil-2-(3-oxobutil)ciclo-hexano-1,3-diona.
\textbf{8.} Cada novo \omterm{def:b2:enolates-aldol:enolate}{enolato} toma um próton de uma molécula de diona (ou do solvente) e assim
regenera o \omterm{def:b2:enolates-aldol:enolate}{enolato} doador: a base não é consumida.
\textbf{9.} O \omterm{def:b2:enolates-aldol:enolate}{enolato} é ambidentado; com um eletrófilo carbonado mole, sob \omterm{def:b2:frontier-orbitals:control}{controle orbital}, ele reage
pelo carbono, e o aduto \ce{C-C} conserva duas ligações \ce{C=O}, a combinação mais forte; um
aduto em \ce{O}, se formado, se reverteria.
\textbf{10.} Um: C2, ligado a C1, C3, à metila e à cadeia.
\textbf{11.} O carbono carbonílico da cadeia e a carbonila do anel C1 (ou C3): \ce{C(=O)}--\ce{CH2}--\ce{CH2}--C2--C1,
cinco carbonos contados inclusive.
\textbf{12.} Para consumir toda a diona, o reagente mais valioso; apenas leve, porque a but-3-en-2-ona
é muito tóxica e polimeriza.
\textbf{13.} C4 e C6 do anel (vizinhos de C3 e C1), o \ce{CH2} da cadeia vizinho de sua carbonila, e o
\ce{CH3} terminal da cadeia. C2 não tem nenhum.
\textbf{14.} O \omterm{def:b2:enolates-aldol:enolate}{enolato} do \ce{CH3} terminal ataca uma carbonila do anel: o anel fechado conta esse
carbono, o carbono carbonílico da cadeia, dois \ce{CH2}, C2 e o carbono carbonílico atacado: seis átomos.
\textbf{15.} O \omterm{def:b2:enolates-aldol:enolate}{enolato} do \ce{CH2} interno da cadeia fecharia um anel de quatro membros, tensionado. Um
\omterm{def:b2:enolates-aldol:enolate}{enolato} do anel (C4 ou C6) atacando a carbonila da cadeia também fecha um anel de seis membros, mas em
ponte: esse \omterm{def:b2:enolates-aldol:aldol}{aldol} não pode desidratar (a ligação dupla ficaria em uma cabeça de ponte de um pequeno
sistema bicíclico) e se reverte, enquanto o \omterm{def:b2:enolates-aldol:aldol}{aldol} fundido é drenado pela desidratação.
\textbf{16.} O \omterm{def:b2:enolates-aldol:aldol}{aldol} leva \ce{OH} no antigo carbono carbonílico do anel, agora na fusão dos anéis; a
perda de água com um hidrogênio do antigo carbono \ce{CH3} (agora em $\alpha$ da cetona da cadeia) dá a
\ce{C=C} conjugada com essa cetona.
\textbf{17.} A nova ligação dupla é conjugada com a carbonila: uma eliminação E1cB passando pelo
\omterm{def:b2:enolates-aldol:enolate}{enolato}, favorecida pelo aquecimento (\cref{prop:b2:enolates-aldol:aldol-equilibrium}).
\textbf{18.} Uma cetona saturada e uma cetona $\alpha,\beta$-insaturada (uma ciclo-hexenona).
\textbf{19.} O carbono de fusão que leva o grupo metila: quatro substituintes diferentes.
\textbf{20.} Não: a diona é aquiral, e suas duas carbonilas são atacadas igualmente (são imagens
especulares uma da outra em relação à cadeia); com reagentes aquirais o produto é racêmico.
Catalisadores aminados quirais favorecem uma carbonila e dão um enantiômero em excesso (volume do
3.º ano).
\textbf{21.} Diona \ce{C7H10O2}: \qty{126.155}{g/mol}; \ce{C4H6O}: \qty{70.091}{g/mol}; \ce{C11H14O2}:
\qty{178.231}{g/mol}.
\textbf{22.} \ce{C7H10O2 + C4H6O -> C11H14O2 + H2O}: C 11, H 16, O 3 de cada lado.
\textbf{23.} $10.0/126.155 = \qty{0.07927}{mol} = \qty{79.3}{mmol}$.
\textbf{24.} $0.07927 \times 178.231 \times 0.70 = \boldsymbol{\approx \qty{9.89}{g}}$ de
cetona de Wieland--Miescher.
\end{solution}
